
\documentclass[sigplan,10pt]{acmart}

\settopmatter{printfolios=false,printacmref=false}
\setcopyright{none}
\acmConference[CSCI 2951X]{CSCI 2951X Course Project}{Spring 2026}{Providence, RI, USA}
\acmYear{2026}
\acmDOI{}
\acmISBN{}
\acmPrice{}

% ── packages ────────────────────────────────────────────────
\usepackage{booktabs}
\usepackage{enumitem}
\usepackage{amsmath}

% ============================================================
\begin{document}

\title{Homeostatic Objectives for Sequential Reinforcement Learning}

\author{Ziyu Kong}
\affiliation{%
  \department{Department of Computer Science}
  \institution{Brown University}
  \city{Providence} \state{RI} \country{USA}}
\email{ziyu_kong@brown.edu}

\author{Shihang Gui}
\affiliation{%
  \department{Department of Computer Science}
  \institution{Brown University}
  \city{Providence} \state{RI} \country{USA}}
\email{shihang\_gui@brown.edu}

\author{Ruth Ukubay}
\affiliation{%
  \department{Department of Computer Science}
  \institution{Brown University}
  \city{Providence} \state{RI} \country{USA}}
\email{ruth\_ukubay@brown.edu}

\author{Meiyi Song}
\affiliation{%
  \department{Department of Computer Science}
  \institution{Brown University}
  \city{Providence} \state{RI} \country{USA}}
\email{meiyi\_song1@brown.edu}







\begin{abstract}
Task-specific reward shaping is effective but brittle. It must be
redesigned whenever the task changes, and the agent cannot reuse what
it learned. This project tests whether a homeostatic objective,
one that never changes and requires no task-specific engineering,
can serve as a reusable control principle across a sequence of related
reinforcement learning tasks. We compare a homeostatic agent, which
optimizes drive reduction from an internal energy setpoint, against a
task-specific agent across a five-task sequential gridworld benchmark
where all tasks share the same embodied energy dynamics. Our pilot
results suggest a consistent tradeoff where the task-specific agent learns earlier tasks faster, while the homeostatic agent transfers more
effectively to later tasks that depend heavily on resource management
and hazard avoidance, and ends successful episodes with higher
remaining energy. We evaluate this tradeoff systematically using
forward transfer, forgetting, and safety margin metrics across
multiple seeds.
\end{abstract}
 
\maketitle
 
% \section{Introduction}
% \lipsum[2-3]
 
\section{Related Work}

Standard RL agents need a new reward function every time the task
changes. Our work sits at the intersection of two areas that try to
address this: homeostatic RL, which grounds reward in the agent's
internal state, and transfer RL, which studies how agents reuse
knowledge across tasks.

\subsection{Homeostatic Reinforcement Learning}

Keramati and Gutkin~\cite{keramati2014} formalize homeostatic RL by
defining reward as the reduction of deviation from an internal setpoint.
The agent is not told what is good or bad; it learns from how actions
affect its own internal state. Their work shows this produces sensible
behaviors like risk aversion and anticipatory regulation, but only
studies a single, unchanging task.

Yoshida~\cite{yoshida2018} shows that richer survival behaviors
emerge when this objective is scaled to deep RL. More recently,
Yoshida~\cite{Yoshida} argues that an HRRL agent with recurrent
connections can adapt when its environment shifts unexpectedly, drawing
a connection to meta-RL. This is the closest prior work to ours, but it
still studies one environment with one changing variable. It does not
ask what happens when the task itself changes.

\subsection{Transfer Across Tasks}

Barreto et al.~\cite{Barreto} propose a framework where an agent
learns a representation of environment dynamics, called successor
features, that can be reused when the reward function changes. When
a new task arrives, the agent combines these stored features with a
new task description using a procedure called generalized policy
improvement, which picks the best available behavior without learning
from scratch. This works well in practice but requires the designer to
specify or learn a feature set that describes all future tasks in
advance. There is no internal state involved; transfer happens
through explicit structural knowledge about the reward.

Tomov et al.~\cite{Tomov} show that humans solving multi-task
problems behave like this too: they reuse old policies rather than
replanning from scratch. Importantly, they note that the task
descriptions their framework requires could be grounded in internal
physiological states like hunger or fatigue,
exactly the connection our project tests.

\subsection{What is Missing}

No existing work tests whether a homeostatic objective alone is enough
to transfer across a sequence of changing tasks in reinforcement
learning. Keramati and Gutkin study one task. Yoshida studies one
shifting environment. Barreto transfers across tasks but needs
task-specific feature engineering and has no internal state. We ask
whether a fixed homeostatic objective, one that never changes and
requires no reward redesign, can serve as a reusable control
principle across a sequential RL task family, precisely because the
tasks share the same underlying embodied resource constraints. If so,
this suggests that homeostasis is not just a reward trick for a single
task, but a structural property that makes RL agents more adaptable
when the world changes around them.

\section{Motivation and Research Question}
\subsection{Motivation}
Conventional reinforcement learning usually depends on externally
specified rewards that are tailored to a particular task. This works
well when the task is fixed, but it becomes brittle in sequential
settings where the objective changes over time. Each task switch may
require redesigned reward shaping, which limits the agent's ability to
reuse previously learned behavior~\cite{Yoshida2016}. This is exactly
the setting in which continual or lifelong learning becomes important:
an autonomous agent should keep adapting in a non-stationary world
without repeatedly changing its objective~\cite{Parisi2019}.

Homeostatic regulation offers a different design principle. Instead of
defining value through task-specific external rewards, it grounds
behavior in the regulation of internal variables toward preferred
setpoints~\cite{keramati2014,konidaris2006}. This is especially
appealing when tasks differ in their immediate goals but share the same
underlying embodied constraints. In our benchmark, every task uses the
same energy dynamics, while the behavioral demands differ. A fixed
homeostatic objective could therefore encourage reusable strategies such
as conserving energy, seeking resources, and avoiding hazards, rather
than forcing the agent to relearn them from scratch after each task
change.
\subsection{Research Question}
We investigate whether a fixed homeostatic objective can improve
sequential reinforcement learning when tasks share embodied
constraints. In the current pilot, we compare a homeostatic agent
against a task-specific agent on a five-task sequential gridworld. In
the full study, we will add standard, energy-aware, and task-specific
oracle baselines. The central question is whether homeostatic
regulation yields better forward transfer, lower forgetting, and
stronger safety margins, even if task-specific shaping learns the
earliest tasks faster.

 
\section{Evaluation Plan}
\subsection{Success rate matrix}
In our sequential setting, an agent is trained on a task sequence
$T_1, T_2, \ldots, T_n$. After finishing training on task $T_i$, we
evaluate the agent on every task seen so far and record the success
rate, defined as the fraction of evaluation episodes in which the task
is completed. This produces an evaluation matrix $R$, where each entry
$r_{i,j}$ denotes the success rate on task $T_j$ after training through
task $T_i$. This style of evaluation follows standard continual RL
practice, but we use a simplified version of the metrics to keep the
analysis interpretable~\cite{continualworld2021,cora2022}. The matrix
below illustrates the setup:
\begin{table}[h]
\centering
\begin{tabular}{lccc}
\hline
\textbf{Train Step} & \textbf{T1} & \textbf{T2} & \textbf{T3} \\
\hline
After T1 & $r_{1,1}$ & --   & --   \\
After T2 & $r_{2,1}$ & $r_{2,2}$ & --   \\
After T3 & $r_{3,1}$ & $r_{3,2}$ & $r_{3,3}$ \\
\hline
\end{tabular}
\caption{Evaluation matrix for sequential task learning.}
\label{tab:sequential_success}
\end{table}

Based on this matrix, we focus on three standard aspects of sequential
learning: forgetting, forward transfer, and backward transfer.
\subsection{Forgetting}
For an earlier task $T_j$, forgetting is the drop from its best
post-training performance to its final performance at the end of the
sequence:
\[
F_j = \max_{i \in \{j, \ldots, n\}} r_{i,j} - r_{n,j}.
\]
Larger values indicate stronger forgetting.
\subsection{Forward transfer} 
Forward transfer asks whether earlier tasks make a new task easier to
learn. For each task $T_i$, we compare the sequential agent's learning
curve on $T_i$ to a single-task baseline trained on $T_i$ alone.
Positive forward transfer means that prior experience helps the agent
adapt to the new task more quickly.
\subsection{Backward transfer}
Backward transfer asks whether learning later tasks changes performance
on earlier ones. In our setting, we summarize it as the change in
success rate on a previously learned task before and after later
training, for example
\[
B_{i,j} = r_{i,j} - r_{j,j}, \qquad i > j.
\]
Positive values indicate that later training improved earlier behavior,
while negative values indicate interference.
\subsection{Evaluation expectation}
A strong sequential agent in our benchmark should exhibit:
\begin{itemize}
    \item high success on the current task,
    \item low forgetting on earlier tasks,
    \item positive forward transfer after task switches, and
    \item stronger safety margins, measured by remaining energy and hazard avoidance.
\end{itemize}

\section{Work Plan}

We organize the project into four stages: agent specification, benchmark finalization, empirical evaluation, and writing. The main goal is to determine whether homeostatic objectives provide a reusable control principle across related sequential tasks, rather than only improving performance on a single handcrafted task.

\subsection{Stage 1: Agent Specification}
The first stage is to finalize the set of agent definitions used throughout the project. This is important because the interpretation of the results depends on whether an agent merely observes internal state or actively optimizes it. We therefore plan to distinguish several agent types in the full study: a standard task-reward agent, an energy-aware task-reward agent, a simplified homeostatic agent with setpoint-based drive reduction, and a task-specific oracle agent trained independently for each task. During this stage, we will ensure that the architecture and training algorithm are matched across agents, so that the main differences come from observation and reward design rather than implementation details.

\subsection{Stage 2: Benchmark Finalization}
In the second stage, we will finalize the sequential benchmark itself. The benchmark will consist of a family of related gridworld tasks that share the same action space, transition dynamics, energy variable, and failure condition, while varying in the degree to which successful behavior depends on resource management and hazard avoidance. We will refine the existing five-task sequence and, if time permits, add control sequences with weaker dependence on shared energy constraints. The deliverable for this stage is a fixed benchmark specification and stable task definitions.

\subsection{Stage 3: Empirical Evaluation}
In the third stage, we will run the main multi-seed experiments and
compare the agents on sequential learning performance. The evaluation
will focus on current-task success, adaptation speed after task
switches, retention on earlier tasks, remaining energy at the end of
successful episodes, and hazard contacts. The main comparison will test
whether the homeostatic agent trades slower early-task learning for
better transfer and stronger safety margins on later related tasks. We
will also compare the sequential agents against task-specific oracle
agents in order to characterize the tradeoff between reusability and
per-task specialization. More broadly, this stage follows continual RL
evaluation frameworks that emphasize transfer, retention, and
post-switch adaptation rather than only final return~\cite{continualworld2021,cora2022}.

\subsection{Stage 4: Analysis and Writing}
In the final stage, we will analyze the results, connect them to prior work on homeostatic reinforcement learning and sequential/continual RL, and consolidate the findings into the final written proposal and report. The expected outcome is a clear argument about where homeostatic objectives are beneficial: not necessarily in maximizing single-task reward, but in supporting reusable and safer behavior across related task sequences. We will also document limitations, including simplified internal-state modeling and the gap between the current benchmark and richer environments such as MiniHack.

\begin{table}[t]
\centering
\small
\begin{tabular}{lcccc}
\toprule
Agent & Task R. & Energy Obs. & Homeo R. & Seq. \\
\midrule
Standard-Seq.      & Y & N & N & Y \\
Energy-Aware-Seq.  & Y & Y & N & Y \\
Homeostatic-Seq.   & Min. & Y & Y & Y \\
Task Oracle        & Y & Opt. & N & N \\
\bottomrule
\end{tabular}
\caption{Planned agent definitions for the full benchmark. ``Min.'' means sparse task success/failure signals are retained, while the main optimization pressure comes from homeostatic reward. The current pilot includes only the task-specific and homeostatic agents.}
\label{tab:agent_types}
\end{table}

\section{Pilot Findings}

\subsection{Pilot Setup}
As an initial validation step, we constructed a simple sequential gridworld benchmark to test whether homeostatic objectives may provide an advantage in a sequence of related reinforcement learning tasks. All tasks shared the same embodied energy dynamics, action space, and transition rules. In every task, the agent moved on a grid, consumed energy at each step, and failed if its energy dropped to zero. The tasks differed only in how strongly successful behavior depended on resource management and hazard avoidance. The current pilot sequence consists of five tasks: \textit{Reach}, where the agent can directly navigate to the goal; \textit{Recharge}, where it must first collect food to survive; \textit{Hazard Reach}, where it must avoid harmful cells; \textit{Detour}, where it must both recharge and avoid hazards; and \textit{Tight Detour}, a stricter version with less margin for error.

We compared two agents using the same DQN architecture and training
protocol. In this pilot phase, we implemented only the task-specific
and homeostatic agents, leaving the energy-aware and task-specific
oracle baselines for the full benchmark. The first agent was a
\textit{task-specific agent}, which received reward shaping tailored to
the active task, including progress-based reward and task-relevant
bonuses or penalties. The second was a \textit{homeostatic agent},
which optimized a shared internal objective based on one-step reduction
in distance to an energy setpoint, together with minimal success and
failure signals. This simplified drive-reduction design is directly
inspired by prior work on homeostatic reinforcement learning and
adaptive motivational systems~\cite{keramati2014,yoshida2018,konidaris2006}.
Both agents were trained sequentially across the five tasks without
resetting model parameters between tasks. We evaluated them using
current-task success, adaptation speed after task switches, retention on
earlier tasks, remaining energy at the end of evaluation episodes, and
hazard contacts.

\subsection{Preliminary Results}
The pilot results suggest a consistent tradeoff rather than a universal dominance of one agent over the other. The task-specific agent learned the earliest and simplest tasks more quickly, which is expected because its reward directly encoded the task objective. However, as the task sequence became more dependent on shared resource-management and safety structure, the homeostatic agent became increasingly competitive and in several cases more robust. In later tasks such as \textit{Hazard Reach}, \textit{Detour}, and \textit{Tight Detour}, the homeostatic agent matched or exceeded the task-specific agent in success while often maintaining a higher remaining energy level, indicating safer or more conservative behavior.

\begin{figure}[t]
    \centering
    \includegraphics[width=\columnwidth]{figures/sequential_homeostasis.png}
    \caption{Pilot results on the five-task sequential gridworld benchmark. Left: current-task success over sequential training. Right: final retention across all tasks after training on the full sequence. The task-specific agent learns the earliest tasks faster, while the homeostatic agent shows stronger retention and more robust performance on later resource- and hazard-dependent tasks.}
    \label{fig:pilot_sequential_homeostasis}
\end{figure}

At the same time, the pilot also shows that both agents still suffer from forgetting, especially on earlier tasks after later sequential training. This means the current evidence does not yet support the stronger claim that homeostatic objectives solve continual reinforcement learning in general. Instead, the current conclusion is narrower and more defensible: homeostatic objectives appear to trade off slower early single-task learning for better transfer, stronger safety margins, and more reusable behavior in later tasks that share the same embodied resource constraints. These pilot findings motivate a larger and more systematic evaluation, rather than serving as final proof on their own.



 
\bibliographystyle{ACM-Reference-Format}
\begin{thebibliography}{99}
 
\bibitem{keramati2014}
Mehdi Keramati and Boris Gutkin. Homeostatic reinforcement learning for integrating reward collection and physiological stability. \emph{eLife}, 3:e04811, 2014.

\bibitem{yoshida2018}
Naoto Yoshida. Homeostatic Agent for General Environment. \emph{Journal of Artificial General Intelligence}, 8(1):1--22, 2018.

\bibitem{Yoshida}
Naoto Yoshida. Meta-Reinforcement Learning in Homeostatic Regulation. Extended abstract, Computational and Cognitive Neuroscience, 2025.

\bibitem{Barreto}
Andre Barreto, Will Dabney, Remi Munos, Jonathan J. Hunt, Tom Schaul, Hado P. van Hasselt, and David Silver. Successor Features for Transfer in Reinforcement Learning. In \emph{Advances in Neural Information Processing Systems 30}, 2017.

\bibitem{Tomov}
Momchil S. Tomov, Eric Schulz, and Samuel J. Gershman. Multi-task reinforcement learning in humans. \emph{Nature Human Behaviour}, 5(6):764--773, 2021.

\bibitem{Yoshida2016}
Naoto Yoshida. On Reward Function for Survival. \emph{arXiv preprint arXiv:1606.05767}, 2016.

\bibitem{Parisi2019}
German I. Parisi, Ronald Kemker, Jose L. Part, Christopher Kanan, and Stefan Wermter. Continual Lifelong Learning with Neural Networks: A Review. \emph{Neural Networks}, 113:54--71, 2019.

\bibitem{konidaris2006}
George Konidaris and Andrew Barto. An Adaptive Robot Motivational System. In \emph{Proceedings of the 9th International Conference on the Simulation of Adaptive Behavior}, pages 346--356, 2006.

\bibitem{continualworld2021}
Michal Wolczyk, Razvan Pascanu, Lukasz Zajac, Andrei P. Rusu, and Piotr Milos. Continual World: A Robotic Benchmark for Continual Reinforcement Learning. In \emph{Advances in Neural Information Processing Systems 34}, 2021.

\bibitem{cora2022}
Sam Powers, Eliot Xing, Eric Kolve, Roozbeh Mottaghi, and Abhinav Gupta. CORA: Benchmarks, Baselines, and Metrics as a Platform for Continual Reinforcement Learning Agents. In \emph{Proceedings of the 1st Conference on Lifelong Learning Agents}, 2022.
\end{thebibliography}
 
\end{document}

 
